\cleardoublepage % memastikan bab baru dimulai di halaman ganjil (kanan)
\chapter{PENDAHULUAN}
\label{chap:pendahuluan}
% --- Latar Belakang ---
\section{Latar Belakang}
Perkembangan teknologi digital telah membawa perubahan fundamental dalam cara negara membangun, mengelola, dan mempertahankan sistem pertahanannya.
Lingkungan strategis global saat ini ditandai oleh meningkatnya kompleksitas ancaman yang bersifat multidimensi, tidak hanya dalam bentuk ancaman militer konvensional, tetapi juga ancaman siber, disrupsi informasi, serta perang berbasis data dan teknologi.
Dalam konteks tersebut, keunggulan pertahanan tidak lagi semata ditentukan oleh kekuatan alutsista atau jumlah personel, melainkan oleh kemampuan mengelola informasi, mengintegrasikan sistem, serta mengambil keputusan secara cepat dan akurat berbasis data.

Transformasi digital menjadi agenda strategis bagi banyak negara dalam memperkuat kapabilitas pertahanan modern.
Negara-negara maju telah menjadikan teknologi digital, data, dan sistem terintegrasi sebagai fondasi utama dalam membangun smart defence, network-centric warfare, dan joint all-domain operations.
Pendekatan ini menempatkan data sebagai aset strategis dan menjadikan interoperabilitas sistem lintas domain sebagai kunci keunggulan operasional.
Dengan demikian, transformasi digital di sektor pertahanan bukan sekadar proses digitalisasi administrasi, melainkan perubahan menyeluruh pada tata kelola, proses bisnis, arsitektur sistem, dan budaya organisasi pertahanan.

Di tingkat nasional, Pemerintah Republik Indonesia telah menetapkan transformasi digital sebagai agenda pembangunan jangka panjang melalui berbagai regulasi dan kebijakan strategis, seperti Sistem Pemerintahan Berbasis Elektronik (SPBE), Percepatan Transformasi Digital Nasional, serta Rencana Pembangunan Jangka Panjang Nasional (RPJPN) 2025–2045.
Kebijakan tersebut menegaskan bahwa seluruh kementerian dan lembaga, termasuk Kementerian Pertahanan Republik Indonesia (Kemhan RI), dituntut untuk membangun ekosistem digital yang terintegrasi, aman, dan berkelanjutan guna mendukung penyelenggaraan pemerintahan dan pembangunan nasional, termasuk dalam sektor pertahanan.

Namun demikian, hingga saat ini implementasi transformasi digital di lingkungan Kemhan RI masih menghadapi berbagai tantangan mendasar.
Sistem informasi dan aplikasi pertahanan masih berkembang secara terpisah di berbagai unit kerja dan matra, sehingga menimbulkan fragmentasi sistem, redundansi aplikasi, serta lemahnya integrasi dan interoperabilitas data.
Kondisi ini berdampak pada terbatasnya kemampuan analisis lintas domain, lambatnya proses pengambilan keputusan strategis, serta belum optimalnya pemanfaatan data sebagai aset pertahanan.
Selain itu, meningkatnya ancaman siber terhadap infrastruktur dan data pertahanan menuntut adanya fondasi arsitektur digital yang lebih kuat, terstandar, dan berorientasi pada keamanan sejak tahap perancangan.

Di sisi lain, Kemhan RI juga dihadapkan pada tuntutan untuk mewujudkan modernisasi pertahanan yang selaras dengan konsep pertahanan cerdas (smart defence).
Konsep ini menekankan pentingnya integrasi antara manusia, proses, teknologi, dan kebijakan dalam satu arsitektur pertahanan digital yang adaptif dan resilien.
Tanpa adanya kerangka arsitektur transformasi digital yang jelas dan terstruktur, upaya modernisasi pertahanan berisiko berjalan parsial, tidak terarah, dan sulit diimplementasikan secara berkelanjutan.

Berdasarkan kondisi tersebut, diperlukan suatu pendekatan sistematis untuk merancang arsitektur transformasi digital yang dapat menjadi arah strategis bagi Kemhan RI.
Arsitektur ini diharapkan mampu mengintegrasikan aspek bisnis, data, aplikasi, dan teknologi pertahanan secara holistik, sekaligus selaras dengan kebijakan nasional dan praktik terbaik internasional.
Penggunaan enterprise architecture framework seperti TOGAF menjadi relevan karena menyediakan metodologi terstruktur untuk menyelaraskan kebutuhan strategis organisasi dengan solusi teknologi secara menyeluruh.

Oleh karena itu, penelitian ini difokuskan pada perancangan arsitektur transformasi digital pada Kementerian Pertahanan Republik Indonesia sebagai landasan dalam mewujudkan pertahanan cerdas yang modern, terintegrasi, dan berdaya tahan tinggi terhadap dinamika ancaman di era digital.
Hasil penelitian diharapkan dapat memberikan kontribusi strategis, baik secara akademik maupun praktis, dalam mendukung pembangunan ekosistem digital pertahanan nasional yang berkelanjutan.

% --- Rumusan Masalah ---
\section{Rumusan Masalah}
Berdasarkan latar belakang yang telah diuraikan, terdapat sejumlah permasalahan mendasar yang menghambat percepatan transformasi digital di lingkungan Kementerian Pertahanan Republik Indonesia.
Permasalahan tersebut muncul sebagai konsekuensi dari fragmentasi sistem, belum optimalnya tata kelola data dan teknologi, serta belum adanya kerangka arsitektur transformasi digital yang dapat menjadi acuan terintegrasi.
Sampai dengan saat ini sistem informasi pertahanan masih berjalan secara terpisah antar unit dan matra, sehingga menghambat integrasi data dan kolaborasi lintas organisasi.
Selain itu, kemajuan teknologi dan meningkatnya ancaman siber menuntut Kemhan untuk membangun fondasi digital yang lebih kuat, terarah, dan sesuai dengan mandat regulatif nasional.

Untuk menjawab tantangan tersebut, penelitian ini merumuskan beberapa pertanyaan pokok sebagai berikut:
\begin{enumerate}
\item	Bagaimana kondisi eksisting (as-is) teknologi, proses bisnis, data, dan sumber daya manusia di lingkungan Kementerian Pertahanan yang berkaitan dengan transformasi digital?
\item	Bagaimana merancang arsitektur transformasi digital Kemhan RI yang relevan dengan kebutuhan pertahanan modern?
\end{enumerate}
Rumusan masalah tersebut menjadi dasar bagi penelitian ini dalam menghasilkan arsitektur transformasi digital yang dapat digunakan Kemhan RI sebagai pedoman strategis dalam membangun ekosistem digital pertahanan yang modern, aman, dan berkelanjutan.

% --- Tujuan ---
\section{Tujuan Penelitian}
Penelitian ini bertujuan untuk merancang suatu arsitektur transformasi digital yang komprehensif bagi Kementerian Pertahanan Republik Indonesia sebagai landasan strategis dalam memperkuat tata kelola pertahanan modern.
Upaya ini dilakukan untuk menjawab tantangan fragmentasi sistem, lemahnya integrasi data pertahanan, meningkatnya ancaman siber, serta tuntutan regulatif nasional yang mewajibkan Kemhan memiliki arsitektur dan peta rencana SPBE yang terstruktur dan terintegrasi.
Secara khusus, penelitian ini memiliki tujuan sebagai berikut:
\begin{enumerate}
    \item Mengidentifikasi kondisi eksisting (as-is) ekosistem digital di lingkungan Kementerian Pertahanan, mencakup proses bisnis, aliran data, aplikasi, infrastruktur teknologi,
    serta tingkat kesiapan sumber daya manusia terhadap transformasi digital.
    \item Merancang arsitektur transformasi digital Kementerian Pertahanan berdasarkan kerangka TOGAF yang mencakup empat domain inti,
    yaitu: arsitektur bisnis, data, aplikasi, dan teknologi.
    \item Memberikan rekomendasi strategis berupa roadmap implementasi transformasi digital Kemhan RI yang berfungsi sebagai
    panduan jangka menengah bagi organisasi dalam mengembangkan dan mengimplementasikan berbagai inisiatif digital secara bertahap, terarah, dan berkelanjutan.
\end{enumerate}
Melalui pencapaian tujuan-tujuan tersebut, penelitian ini diharapkan dapat memberikan kontribusi nyata dalam mendukung modernisasi pertahanan nasional
melalui pembangunan arsitektur digital yang kuat, adaptif, dan selaras dengan kebutuhan keamanan nasional di era digital.

% --- Manfaat Penelitian ---
\section{Manfaat Penelitian}
Penelitian mengenai desain arsitektur transformasi digital pada Kementerian Pertahanan RI diharapkan memberikan kontribusi teoretis 
maupun praktis yang relevan bagi pengembangan ekosistem digital pertahanan nasional.

\subsection{Manfaat Teoritis}
\begin{enumerate}
    \item Menghasilkan desain arsitektur transformasi digital pada Kementerian Pertahanan dengan menggunakan kerangka kerja TOGAF.
    \item Menjadi referensi akademik bagi penelitian lanjutan dalam bidang smart defense, AI governance, dan digital military transformation.
\end{enumerate}

\subsection{Manfaat Praktis}
\begin{enumerate}
    \item Memberikan rekomendasi dan model desain arsitektur transformasi digital bagi Kementerian Pertahanan.
    \item Mendukung peningkatan tata kelola organisasi, integrasi sistem pertahanan dan modernisasi pertahanan.
    \item Memperkuat kapabilitas keamanan siber Kemhan.
\end{enumerate}
Dengan demikian, penelitian ini tidak hanya memberikan nilai ilmiah, tetapi juga memberikan arah strategis yang dapat langsung diterapkan
pada upaya transformasi digital di lingkungan Kementerian Pertahanan RI.

% --- Ruang Lingkup Penelitian ---
\section{Ruang Lingkup Penelitian}
Ruang lingkup penelitian ini ditetapkan untuk memastikan fokus analisis yang jelas dan sesuai dengan tujuan penelitian.
Adapun batasan dan cakupan penelitian ini meliputi:
\begin{enumerate}
    \item Lingkup organisasi: Kementerian Pertahanan Republik Indonesia, khususnya unit-unit kerja yang terkait dengan perencanaan,
    pengelolaan dan implementasi transformasi digital, seperti Ditjen Pothan, Pusdatin Kemhan dan Bainstrahan Kemhan.
    \item Lingkup Teknis: Arsitektur bisnis, data, aplikasi, dan teknologi berdasarkan kerangka kerja TOGAF ADM.
    \item Batasan: Penelitian ini tidak mencakup aspek rahasia militer dan operasi tempur, namun fokus pada integrasi sistem informasi, keamanan data, dan tata kelola digital.
\end{enumerate}

% --- Keluaran Penelitian ---
\section{Keluaran Penelitian}
Penelitian ini diarahkan untuk menghasilkan keluaran yang bersifat konseptual maupun aplikatif dalam mendukung transformasi digital Kementerian Pertahanan RI.
Adapun keluaran yang dihasilkan meliputi:
\begin{enumerate}
    \item Model arsitektur transformasi digital Kemhan RI dengan menggunakan pendekatan TOGAF ADM Phase A–D, mencakup arsitektur bisnis, data, aplikasi, dan teknologi.
    \item Rekomendasi Strategis dan Roadmap Implementasi Transformasi Digital Kemhan 2026–2030.
\end{enumerate}


% % --- Metodologi Pengerjaan TA ---
% \section{Metodologi}
% Tuliskan semua tahapan spesifik yang akan dilalui selama pelaksanaan tesis. 
% Ada beberapa pendekatan metodologi yang dapat digunakan dalam pelaksanaan tesis sesuai dengan topik yang dibahas, 
% misalnya Design Science Research Methodology (DSRM), Design Research Methodology (DRM), Cross-Industry Standard Process for Data Mining (CRISP-DM), dan sebagainya.

% Secara umum,  metodologi pengerjaan tesis meliputi tahapan, langkah-langkah, atau tata cara melakukan:
% \begin{enumerate}
% \item	Investigasi pengumpulan fakta di latar belakang untuk merumuskan masalah.
% \item	Pencarian, pengelompokan, dan penapisan literatur atau sumber informasi untuk mengumpulkan informasi yang relevan tentang topik yang diangkat, termasuk teori (konsep atau teori apa saja yang perlu dicari), hal-hal yang telah dicapai oleh orang lain (cara mencari dan kata kuncinya), dan berbagai informasi pendukung, untuk mencari solusi terhadap masalah yang dibahas. Gunakan metodologi yang tepat dalam menggali informasi dan dokumentasikan prosesnya (termasuk rekaman wawancara atau survei) di dalam Lampiran, termasuk tautan ke video atau foto. Hasil penggalian informasi ini akan dijelaskan secara komprehensif di Bab II Studi Literatur.
% \item   Analisis kebutuhan berdasarkan hasil investigasi dan studi literatur. Jelaskan metode yang digunakan untuk melakukan analisis kebutuhan, misalnya wawancara, survei, observasi, atau studi kasus. Hasil analisis kebutuhan ini akan dijelaskan secara rinci di bab berikutnya.
% \item   Perancangan solusi berdasarkan hasil analisis kebutuhan. Jelaskan metode atau tahapan perancangan yang digunakan, dengan abstraksi yang bertahap dari generik ke spesifik, misalnya dengan pendekatan \textit{C4 model}. Selain itu, gunakan alat bantu penggambaran diagram yang tepat, misalnya menggunakan diagram UML, SySML, ERD, atau yang lainnya. 
% \item   Jika ada implementasi terhadap hasil rancangan, jelaskan teknologi dan alat bantu yang digunakan untuk mengimplementasikan rancangan. 
% \item   Evaluasi terhadap solusi yang telah diimplementasikan. Jelaskan metode pengujian, metrik evaluasi, atau studi kasus yang digunakan untuk menilai keberhasilan solusi. 
% \item   Penarikan kesimpulan dan saran berdasarkan hasil pelaksanaan tesis berdasarkan hasil evaluasi.   
% \vfill % memaksa konten di atas untuk berada di bagian atas halaman dan menyisakan ruang kosong di bawah konten ini akibat berakhirnya section ini
% \end{enumerate}


% --- Sistematika Penulisan ---
\section{Sistematika Penulisan}
Sistematika penulisan dalam proposal penelitian ini disusun untuk memberikan alur pembahasan yang runtut, logis, dan mudah dipahami.
Setiap bab dirancang untuk saling mendukung satu sama lain, mulai dari pengenalan masalah hingga penyusunan metodologi penelitian yang digunakan dalam merancang arsitektur transformasi digital pada Kementerian Pertahanan RI.
Adapun sistematika penulisan proposal ini adalah sebagai berikut:
\begin{enumerate}
\item	Bab I Pendahuluan: Bab ini menjelaskan landasan awal penelitian, meliputi latar belakang permasalahan, rumusan masalah, tujuan penelitian, manfaat penelitian, ruang lingkup penelitian, keluaran penelitian, serta sistematika penulisan.
Bab ini berfungsi untuk memberikan gambaran umum mengenai konteks, arah, dan batasan penelitian.       .
\item	Bab II Studi Literatur: Bab ini menyajikan kajian teoretis dan konseptual yang menjadi dasar penelitian.
Pembahasan mencakup konsep transformasi digital, transformasi digital di sektor pertahanan, arah dan tantangan transformasi digital Kementerian Pertahanan RI,
serta konsep pertahanan cerdas (smart defence). Selain itu, bab ini juga mengulas framework TOGAF sebagai kerangka kerja arsitektur enterprise yang digunakan dalam perancangan arsitektur transformasi digital.
Tinjauan pustaka ini berfungsi sebagai pijakan teoritis dan konseptual dalam pengembangan artefak penelitian.
\item	Bab III Analisis dan Desain: Bab ini menjelaskan pendekatan penelitian yang digunakan, yaitu Design Science Research Methodology (DSRM),
beserta tahapan-tahapannya mulai dari identifikasi masalah hingga komunikasi hasil penelitian.
Bab ini juga menguraikan jenis dan sumber data, serta metode pengumpulan data yang digunakan untuk menghasilkan artefak arsitektur yang valid dan relevan.
\item   Daftar Pustaka: Berisi daftar referensi yang digunakan dalam penyusunan laporan tesis.
\end{enumerate}